Recent works have demonstrated the utility of LLMs for \emph{text steganography}---i.e., for encoding hidden messages within seemingly innocuous text. For steganographic methods to enable safe communication, they ought to be resilient to meaning-preserving changes to stegotexts. Yet, prior methods are universally susceptible to global paraphrasing of the encoded texts due to their reliance on \emph{form}-related features: since paraphrasing tends to alter such features, the encoded message can easily be scrambled. We address this limitation by introducing a symmetric key method that operates on clearly delimited \emph{semantic} units, yielding greater robustness to paraphrasing. We implement this method for three NLP tasks, validating its robustness under several paraphrase-based attack types. Finally, we present steganalysis results illustrating the difficulty of the distinguishing the resulting stegotexts from covertexts on the same task.\footnote{Code and data will be released upon paper acceptance.}

% To address this limitation, we show that the steganographic protocol from \citet{perry2025robust} may be adapted to operate on clearly delimited \emph{semantic} units, yielding greater robustness to paraphrasing. We empirically validate this robustness by implementing these adaptations for three diverse generation tasks: story generation, summarization, and recipe generation.